% answers_short.tex —— 2015 年·北京卷·文科 · 简短答案速查
% 由 tools/make_short_answers.py 自 answers.tex 的 \gksolhead/\gkanswer 提取，
% 并与 2015/2015_数学真题_北京文_答案要点.md 逐题比对后生成，勿手改（改源 MD 后重跑）。

\begin{gknotice}
\textbf{使用说明}：本节只给最终答案便于快速核对，完整推导见解析版。答案为编写者独立推导，\textbf{非官方标准答案}。
\end{gknotice}

\gkpart{选择题}
\begin{center}\begin{tabular}{ccccccccc}
\toprule
题号 & 1 & 2 & 3 & 4 & 5 & 6 & 7 & 8 \\
\midrule
答案 & A & D & B & C & B & A & C & B \\
\bottomrule
\end{tabular}\end{center}

\gkpart{填空题}
\begin{itemize}[leftmargin=2.6em]
  \item[\textbf{9.}] \(-1\)
  \item[\textbf{10.}] \(\log_{2}5\)
  \item[\textbf{11.}] \(\dfrac{\pi}{4}\)
  \item[\textbf{12.}] \(\sqrt{3}\)
  \item[\textbf{13.}] \(7\)
  \item[\textbf{14.}] ① 乙；② 数学
\end{itemize}

\gkpart{解答题}
\begin{itemize}[leftmargin=2.6em]
  \item[\textbf{15.}] (1) \(2\pi\)；(2) \(-\sqrt{3}\)（在 \(x=\dfrac{2\pi}{3}\) 处取得）
  \item[\textbf{16.}] (1) \(a_{n}=2n+2\)（\(n\in\mathbf{N}^{*}\)）；(2) \(b_{6}\) 与数列 \(\{a_{n}\}\) 的第 \(63\) 项相等（\(b_{6}=a_{63}=128\)）
  \item[\textbf{17.}] (1) \(\dfrac{1}{5}\)；(2) \(\dfrac{3}{10}\)；(3) 丙（同时购买丙的可能性最大，条件概率 \(\dfrac{600}{685}\approx 0.876\)）
  \item[\textbf{18.}] (1) 成立（\(O\)、\(M\) 为中点故 \(OM\parallel VB\)，\(OM\subset\) 平面 \(MOC\)，故 \(VB\parallel\) 平面 \(MOC\)）；(2) 成立（\(VO\perp AB\) 且面面垂直得 \(VO\perp\) 平面 \(ABC\)；\(VO\perp OC\)、\(OC\perp AB\)，故 \(OC\perp\) 平面 \(VAB\)，从而平面 \(MOC\perp\) 平面 \(VAB\)）；(3) \(\dfrac{\sqrt{3}}{3}\)
  \item[\textbf{19.}] (1) 单调递减区间 \(\left(0,\sqrt{k}\right)\)，单调递增区间 \(\left(\sqrt{k},+\infty\right)\)；极小值 \(f(\sqrt{k})=\dfrac{k}{2}(1-\ln k)\)（无极大值）；(2) 成立（存在零点 当且仅当 \(k\geqslant\mathrm{e}\)；此时 \(f\) 在 \((1,\sqrt{\mathrm{e}}]\) 上递减、端点值 \(f(1)=\dfrac{1}{2}>0\geqslant f(\sqrt{\mathrm{e}})=\dfrac{\mathrm{e}-k}{2}\)，故恰有一个零点）
  \item[\textbf{20.}] (1) \(\dfrac{\sqrt{6}}{3}\)；(2) \(1\)；(3) \(BM\parallel DE\)（恒等式对任意过 \(D\) 且不过 \(E\) 的直线成立）
\end{itemize}
